\documentclass[11pt,letterpaper]{mystyle}
\usepackage{fancyhdr}
\usepackage{float}

\usepackage[utf8]{inputenc} 
\usepackage[T1]{fontenc}
\usepackage[numbers]{natbib}
\usepackage{adjustbox}
\usepackage{breakurl}    
\usepackage{hyperref}
\usepackage[edges]{forest}
\usepackage{subcaption}
\usepackage{soul}
\usepackage{multirow}
\usepackage[utf8]{inputenc} 
\usepackage{booktabs}       
\usepackage{amsfonts}       
\usepackage{nicefrac}      
\usepackage{microtype}     
\usepackage{xcolor}        
\usepackage{colortbl}
\usepackage{enumitem}
\usepackage{amssymb}
\usepackage{wrapfig}
\usepackage{courier}
\usepackage{bxcoloremoji}
\usepackage{CJKutf8}
\usepackage{tcolorbox}
\usepackage{array}
\usepackage{tabularx}

\newcommand{\github}{\raisebox{-1.5pt}{\includegraphics[height=1.05em]{logo/github-logo.pdf}}}
\newcommand{\cqgcomment}[1]{{\begin{CJK*}{UTF8}{gbsn}\textcolor{red}{[#1 -- \textsc{CQG}]}\end{CJK*}}}
\fancypagestyle{headstyle}{
    \fancyhead[L]{
        \includegraphics[width=80pt]{logo/larg-logo.pdf}
    }
    \fancyhead[C]{}
    \fancyhead[R]{}

}

\definecolor{hidden-red}{RGB}{205, 44, 36}
\definecolor{hidden-blue}{RGB}{194,232,247}
\definecolor{hidden-orange}{RGB}{243,202,120}
\definecolor{hidden-green}{RGB}{34,139,34}
\definecolor{hidden-pink}{RGB}{255,245,247}
\definecolor{hidden-black}{RGB}{20,68,106}
\definecolor{purple}{RGB}{144,153,196}
\definecolor{yellow}{RGB}{255,228,123}
\definecolor{hidden-yellow}{RGB}{255,248,203}
\definecolor{tkcolor}{RGB}{224,223,255}
\definecolor{darkblue}{rgb}{0, 0.40, 0.75}
\newcommand{\eg}{\textit{e.g.,}}
\hypersetup{colorlinks=true, citecolor=darkblue, linkcolor=darkblue, urlcolor=darkblue}

\newtcolorbox{socialpost}{
  colback=black!5!white,
  colframe=black!75!black,
  sharp corners,
  boxrule=0.5pt,
  fontupper=\sffamily\small, 
  left=3mm,
  right=3mm,
  top=3mm,
  bottom=3mm,
  boxsep=1mm
}


\tcbset{
  aibox/.style={
    width=\linewidth,
    top=8pt,
    bottom=4pt,
    colback=blue!6!white,
    colframe=black,
    colbacktitle=black,
    enhanced,
    center,
    attach boxed title to top left={yshift=-0.1in,xshift=0.15in},
    boxed title style={boxrule=0pt,colframe=white,},
  }
}



\newtcolorbox{AIbox}[2][]{aibox,title=#2,#1}

\tcbset{
  takeawaybox/.style={
    width=\linewidth,
    top=8pt,
    bottom=4pt,
    colback=hidden-yellow,
    colframe=black,
    colbacktitle=black,
    enhanced,
    center,
    attach boxed title to top left={yshift=-0.1in,xshift=0.15in},
    boxed title style={boxrule=0pt,colframe=white,},
  }
}

\newtcolorbox{TakeawayBox}[2][]{takeawaybox,title=#2,#1}

\title{A Survey of Reasoning Hallucinations}

\author{
  Jiannan Guan$^{1*}$ \quad Qiguang Chen$^{1*}$ \quad Linfan Dai$^{2*}$ \quad Libo Qin$^{2, \coloremojicode{2709}}$ \quad Jingxiang Weng$^{1}$ \quad Dengyun Peng$^{1}$ \quad Jinhao Liu$^{1*}$ \quad Zheng Yan$^{1}$ \quad Bofei Liu$^{1}$ \quad Han He$^{1}$ \quad Jiaqi Wang$^{3}$ \quad Mengkang Hu$^{4}$ \vspace{-5pt}\\
  \textbf{Zhi Chen$^{5}$ \quad Yantao Du$^{5}$ \quad Wanxiang Che$^{1, \coloremojicode{2709}}$}\vspace{-5pt}\\
\normalfont{$^1$ LARG, Research Center for Social Computing and Interactive Robotics, Harbin Institute of Technology,\vspace{-5pt}\\
$^2$ School of Computer Science and Engineering, Central South University,\vspace{-5pt}\\
$^3$ The Chinese University of Hong Kong \quad
$^4$ The University of Hong Kong \quad
$^5$ ByteDance China (Seed)\\
}}


\begin{document}

\begin{abstract}
  \vspace{5mm}
  \textbf{\large Abstract:}
  \vspace{2mm}

  As the volume of peer-reviewed research surges, scholars increasingly rely on social platforms for discovery, while authors invest considerable effort in promoting their work to ensure visibility and citations. To streamline this process and reduce the reliance on human effort, we introduce Automatic Promotion (AutoPR), a novel task that transforms research papers into accurate, engaging, and timely public content. To enable rigorous evaluation, we release PRBench, a multimodal benchmark that links 512 peer-reviewed articles to high-quality promotional posts, assessing systems along three axes: Fidelity (accuracy and tone), Engagement (audience targeting and appeal), and Alignment (timing and channel optimization). We also introduce PRAgent, a multi-agent framework that automates AutoPR in three stages: content extraction with multimodal preparation, collaborative synthesis for polished outputs, and platform-specific adaptation to optimize norms, tone, and tagging for maximum reach. When compared to direct LLM pipelines on PRBench, PRAgent demonstrates substantial improvements, including a 604\% increase in total watch time, a 438\% rise in likes, and at least a 2.9x boost in overall engagement. Ablation studies show that platform modeling and targeted promotion contribute the most to these gains. Our results position AutoPR as a tractable, measurable research problem and provide a roadmap for scalable, impactful automated scholarly communication.
  \vspace{5mm}

  $^{*}$ \textit{Equal Contribution}

  $^{\coloremojicode{2709}}$ \textit{Corresponding Author}

  \vspace{5mm}

  \coloremojicode{1F4C5} \textbf{Date}: Oct. 01, 2025

  \github{} \textbf{Project Website}: \href{https://yzweak.github.io/autopr.github.io}{https://yzweak.github.io/autopr.github.io}

  \github{} \textbf{Code Repository}: \href{https://github.com/LightChen2333/AutoPR}{https://github.com/LightChen2333/AutoPR}

  \coloremojicode{1F4DA}  \textbf{Demo}: \href{https://huggingface.co/spaces/yzweak/AutoPR}{https://huggingface.co/spaces/yzweak/AutoPR}

  \coloremojicode{1F917}  \textbf{Benchmark}: \href{https://huggingface.co/datasets/yzweak/PRBench}{https://huggingface.co/datasets/yzweak/PRBench}




  \coloremojicode{1F4E7} \textbf{Contact}: \href{mailto:qgchen@ir.hit.edu.cn}{qgchen@ir.hit.edu.cn}, \href{mailto:zyan@ir.hit.edu.cn}{zyan@ir.hit.edu.cn},\href{mailto:car@ir.hit.edu.cn}{car@ir.hit.edu.cn}, \href{mailto:lbqin@csu.edu.cn}{lbqin@csu.edu.cn}


\end{abstract}
\maketitle

\vspace{3mm}
\pagestyle{headstyle}
\thispagestyle{empty}

\begin{CJK}{UTF8}{gbsn}

\vspace{-2mm}\section{Introduction}\vspace{-1mm}

\begin{figure}[t]
    \centering
    \includegraphics[width=0.98\textwidth]{figures/framework.pdf}
    \caption{Overview of our study: Automatic Promotion (AutoPR) task, its benchmark PRBench, and the associated method PRAgent. The details of citation trend analysis are shown in Appendix~\ref{append:citation}.
    }
    \label{fig:framework}
\end{figure}

The rapid development of Large Language Models (LLMs) has pushed the capabilities of artificial intelligence to new heights~\cite{achiam2023gpt,yang2025qwen3,qin2024large}.
In particular, LLMs that enhance complex task-solving abilities through specific training strategies, such as Chain-of-Thought (CoT), are at the forefront of this progress~\cite{wei2022chain,guo2025deepseek,jaech2024openai}.
We refer to these as \textbf{Large Reasoning Models (LRMs)}~\cite{chen2025towards}.
However, this advancement is accompanied by a growing and significant challenge: Hallucination.

Initially, research on hallucination was primarily concentrated on the factual level~\cite{huang2025survey,park2025steer}.
For instance, early surveys classified hallucinations into two main categories:
Factuality Hallucination, where generated content conflicts with verifiable real-world knowledge, and Faithfulness Hallucination, where the output contradicts the provided source text or context.
Although significant, these errors are mainly concerned with the correctness of the output's result.

However, the nature of hallucination has fundamentally shifted as LRMs increasingly perform multi-step reasoning, particularly driven by techniques like CoT.
The CoT strategy was introduced to improve complex problem-solving by guiding models to generate explicit intermediate steps.
While this approach appears to offer a path to more reliable and interpretable AI, deeper investigation reveals a concerning reality.
The explicit reasoning paths, while offering transparency, simultaneously introduce new and subtle vulnerabilities.
The core of the hallucination problem has shifted from \textit{what the model knows} to \textit{how the model thinks}.
This has led to the emergence of a more complex and concealed category: \textbf{Reasoning Hallucination}.

In this paradigm, a model may produce the correct final answer, yet its reasoning process is flawed with logical fallacies, fabrications, or disconnected steps.
Conversely, a model might follow a seemingly coherent reasoning chain only to arrive at an incorrect conclusion~\cite{joshi2024llms,yao2025reasoning}.
Specifically, the reasoning process of these models can be unreliable.
This unreliability manifests not only as logical fallacies but also as a pathological disregard for simple solutions, where models fall into unnecessary complexity and overthinking~\cite{peng2025revisiting,fan2025missing}.
This behavior is perhaps another byproduct of their reasoning-style text pattern matching mechanism.
This transition from focusing on factual errors to reasoning process flaws represents the current research frontier~\cite{jiao2025trustworthy}.
It also highlights the profound paradoxes underlying advanced reasoning capabilities.

Ironically, while CoT can reduce the frequency of factual hallucinations in some cases, it may also obscure key cues for their detection.
When the model generates a long and complex reasoning chain, its reported confidence in the final output increases significantly, and the extended reasoning process itself appears more persuasive.
This makes hallucinations harder to detect, both for automated evaluation toolkits and human reviewers~\cite{cheng2025chain}.
This creates a fundamental tension between capability and evaluability, which we term the \textbf{Reasoning Evaluability Paradox}.
This paradox is that the very tools we use to enhance model reasoning (like CoT) may simultaneously weaken our ability to assess their reliability.

The main contributions of this work are as follows:

\begin{itemize}[leftmargin=16pt, itemsep=0pt, topsep=0pt]
    \item \textbf{A Systematic Taxonomy of Reasoning Hallucination:}
    We propose a systematic taxonomy for \textbf{Reasoning Hallucination}.
    This framework moves beyond traditional, result-based definitions of hallucination to focus on process-based flaws that occur during complex reasoning.
    We subdivide reasoning hallucination into four main categories (\textbf{Logical Fallacies, Flaw Repetition, Think-Answer Mismatch, and Overthinking}) to facilitate detailed analysis.

    \item \textbf{A Comprehensive Review of the Reasoning Paradigm:}
    This paper provides a comprehensive review of the causes, detection methods, and mitigation strategies for reasoning hallucination.
    We place special emphasis on analyzing how techniques like CoT fundamentally alter the form of hallucination, shifting it from incorrect answers to flawed reasoning processes.

    \item \textbf{Critical Analysis and Future Roadmap:}
    We conduct a critical analysis of the core challenges in current research, particularly the \textbf{Reasoning Evaluability Paradox} as defined in this paper.
    Finally, we propose future research directions for building more robust, reliable, and trustworthy large reasoning models.
\end{itemize}


\vspace{-2mm}\section{Formalizing Reasoning Hallucination}\vspace{-1mm}
\label{sec:definition}

To establish a rigorous foundation for our analysis, we introduce a formal definition of reasoning hallucination. Our goal is to move beyond ambiguous descriptions and provide a precise criterion that captures the various process-based flaws discussed in this survey.

\subsection{Preliminary}

Let $Q$ represent the input query or problem posed to the model. The model's complete output consists of a final answer $A$ and a reasoning process $R$. We define the reasoning process $R$ as an ordered sequence of $N$ intermediate steps:
$$R = (r_1, r_2, \dots, r_N)$$

To evaluate the quality of this reasoning process, we introduce two key predicates:

\textbf{Step Validity ($V$):} Let $V(r_i | r_{<i}, Q) \in \{\text{True, False}\}$ be a predicate that evaluates whether the reasoning step $r_i$ is a valid (e.g., logically sound) deduction based on the query $Q$ and all preceding steps $r_{<i} = (r_1, \dots, r_{i-1})$. A value of $\text{False}$ indicates a flaw in the step itself, such as a logical fallacy.

\textbf{Process-Answer Consistency ($C$):} Let $C(R, A) \in \{\text{True, False}\}$ be a predicate that evaluates whether the entire reasoning process $R$ logically entails or supports the model's final answer $A$. A value of $\text{False}$ indicates a mismatch between the reasoning and the conclusion, corresponding to Think-Answer Mismatch category.

\subsection{Definition of Valid Reasoning}

Using these predicates, we first define an ideal or \textbf{Valid Reasoning} process. A process is considered valid if and only if every step is valid and the process as a whole is consistent with the final answer.

\newtheorem{definition}{Definition}
\begin{definition}[Valid Reasoning]
    \label{def:valid_reasoning}
    A reasoning process $R$ for a given query $Q$ and answer $A$ is defined as \textbf{Valid}, denoted $Valid(R, Q, A)$, if and only if:
    $$
        Valid(R, Q, A) \equiv \left( \forall i \in [1, N], V(r_i | r_{<i}, Q) = \text{True} \right) \land \left( C(R, A) = \text{True} \right)
    $$
\end{definition}

\subsection{Definition of Reasoning Hallucination}

Our central thesis is that reasoning hallucination encompasses not only \textit{invalidity} but also \textit{pathological sub-optimality}. A process that is correct but unnecessarily complex or redundant (i.e., Overthinking) is also a form of process-based flaw. We therefore present a unified definition of Reasoning Hallucination ($H_R$) that captures both conditions.

\begin{definition}[Reasoning Hallucination]
    \label{def:reasoning_hallucination}
    A reasoning process $R$ is classified as a \textbf{Reasoning Hallucination} ($H_R$) if it is either \textbf{not Valid}, or if it \textbf{is Valid but contains redundant steps} (i.e., is sub-optimal). Formally:
    $$
        H_R(R, Q, A) \equiv \underbrace{\neg Valid(R, Q, A)}_{\text{Invalidity Flaw}} \quad \lor \quad \underbrace{\left( Valid(R, Q, A) \land \exists R' \subset R \text{ s.t. } Valid(R', Q, A) \right)}_{\text{Sub-optimality Flaw}}
    $$
    where $R' \subset R$ denotes that $R'$ is a strict subsequence of $R$, formed by removing one or more steps from $R$.
\end{definition}


\vspace{-2mm}\section{Types of Reasoning Hallucination}\vspace{-1mm}
\label{sec:taxonomy}

Based on the formalization in Definition~\ref{def:reasoning_hallucination}, we systematically categorize reasoning hallucination ($H_R$) into two primary classes.

\subsection{Invalidity Flaws}
\label{sec:invalidity_flaws}

Invalidity Flaws are violations of Valid Reasoning (Definition~\ref{def:valid_reasoning}).
These flaws manifest when at least one reasoning step $r_i$ is logically invalid (i.e., $V(r_i | r_{<i}, Q) = \text{False}$), or when the complete reasoning process $R$ fails to logically entail the final answer $A$ (i.e., $C(R, A) = \text{False}$).

This category encompasses several distinct failure modes.
\textbf{Logical Fallacies} represent the most direct form of invalidity ($V(r_i) = \text{False}$), where models produce steps that are not logically sound.
This includes causal misattributions, such as the ``illusion of causality'' or the ``post hoc fallacy''~\cite{joshi2024llms}, and a ``Fake Reasoning Bias''~\cite{wang2025towards}, where models favor superficial rhetorical structures over logical soundness.

Beyond single-step failures, \textbf{Think-Answer Mismatch} corresponds to the $C(R, A) = \text{False}$ condition.
In this phenomenon, the model's generated reasoning chain $R$ is disconnected from its final answer $A$, suggesting the CoT serves as a post-hoc rationalization rather than the true derivation process~\cite{yao2025reasoning}.
Notably, even when a model exhibits this type of hallucination and reaches an incorrect conclusion during the reasoning process, it may still produce the correct final answer.
This possibility poses a significant challenge to the identification of reasoning hallucination~\cite{yee2024dissociation}.

A distinct type of invalidity emerges as \textbf{Flaw Repetition}, where the model becomes trapped in a loop of semantically similar but incorrect reasoning steps~\cite{yao2025reasoning}.

\subsection{Sub-optimality Flaws}
\label{sec:suboptimality_flaws}

Sub-optimality Flaws represent a more insidious form of $H_R$.
As defined in Definition~\ref{def:reasoning_hallucination}, these are processes that are technically $Valid$ but are pathological due to redundancy (i.e., $\exists R' \subset R \text{ s.t. } Valid(R', Q, A)$).

The primary manifestation of this category is \textbf{Overthinking}.
This phenomenon is characterized by the generation of verbose and tangential reasoning traces even for simple queries~\cite{cuesta2025lrms}.
More pathologically, this behavior causes models to disregard correct solutions even when explicitly provided, leading them to pursue unnecessarily complex and often erroneous reasoning paths.
This suggests that the model is pathologically matching the form of complex reasoning (i.e., generating long token sequences) rather than executing the function of efficient problem-solving.


\vspace{-2mm}\section{Characteristics of Reasoning Hallucination: Process vs. Result}\vspace{-1mm}
\label{sec:characteristics}

The emergence of reasoning hallucination signifies a paradigm shift in the study of LLM reliability, the locus of concern has shifted from \textit{what the model knows} to \textit{how the model thinks}.
This section contrasts reasoning hallucination with traditional hallucination to highlight its process-centric nature.

\subsection{Traditional Hallucination: Result-Centric}

Traditional hallucination, such as Factuality and Faithfulness Hallucination, is fundamentally result-centric.
Its evaluation is external and binary: a model's final answer $A$ is compared against verifiable real-world knowledge $W$ ($A \not\subset W$) or a provided source context $C$ ($A \not\subset C$).
The defect is one of \textit{knowledge}, where either the model lacks the correct fact or fails to faithfully reproduce the context.

\subsection{Reasoning Hallucination: Process-Centric}

In contrast, reasoning hallucination is fundamentally process-centric.
The primary object of scrutiny is not the correctness of the final answer $A$, but the validity and optimality of the reasoning chain $R = (r_1, \dots, r_N)$ used to derive it, which is a defect of \textit{logic}.
A model may produce a correct answer $A$ via a flawed process $R$, or a seemingly coherent process $R$ may lead to an incorrect answer~\cite{joshi2024llms,yao2025reasoning,cuesta2025lrms}.

\subsection{Core Differentiator: The Evaluability Paradox}

This shift from result to process introduces the core challenge defined in this work: the \textbf{Reasoning Evaluability Paradox}.
While traditional factual errors are at least detectable (e.g., via exact match or external search), reasoning flaws are perniciously covert.
The long, complex, and seemingly confident CoT generated by an LRM is often highly persuasive~\cite{cheng2025chain}.
This persuasiveness, a desired outcome of CoT training, ironically obscures the very flaws it generates, making detection by both human reviewers and automated tools exceptionally difficult.
The mechanism used to enhance capability simultaneously degrades our ability to assess its reliability.


\vspace{-2mm}\section{Evaluationof Reasoning Hallucination}\vspace{-1mm}
\label{sec:evaluation}

Addressing the Reasoning Evaluability Paradox requires specialized evaluation tools that can scrutinize the reasoning process, which in turn necessitates a deep understanding of the mechanisms that generate these process-based flaws.

\subsection{Automated Detection Techniques}

Automated detection aims to identify flaws dynamically.
This includes training watchdog classifiers, which have been shown to detect reasoning-induced hallucinations with high accuracy~\cite{jiang2024large}.
Another avenue involves probing the model's internal confidence signals~\cite{chen2024teaching}.
However, this approach is compromised by the overconfidence phenomenon, where models exhibit high confidence for flawed reasoning~\cite{yao2025reasoning}.  % TODO cite gjn reasoning overconfidence

\subsection{Evaluation Benchmarks}

\subsection{Origins and Mechanisms of Reasoning Hallucination}


\vspace{-2mm}\section{Mitigation Strategies}\vspace{-1mm}
\label{sec:mitigation}
xxx


\vspace{-2mm}\section{Frontiers \& Future Direction}\vspace{-1mm}
\label{sec:future}
xxx


\vspace{-2mm}\section{Related work}\vspace{-1mm}
\label{sec:related_work}
xxx


\vspace{-2mm}\section{Conclusion}\vspace{-1mm}
\label{sec:conclusion}
xxx


\end{CJK}

\bibliographystyle{./refstyle}
\bibliography{ref}
\appendix
\newpage

\begin{center}
\textbf{\LARGE Appendix}
\end{center}

\vspace{-2mm}\section{Citation Trend Analysis Details}\vspace{-1mm}
\label{append:citation}

\end{document}